\documentclass[UTF8,11pt]{ctexart}
\usepackage[a4paper,margin=24mm]{geometry}
\usepackage{amsmath,amssymb,amsthm,mathtools,mathrsfs}
\usepackage[all]{xy}
\usepackage{xurl,hyperref,enumitem}
\usepackage{tikz}
\usetikzlibrary{arrows.meta,cd,decorations.markings,calc,patterns}
\hypersetup{hidelinks,breaklinks=true}
\setlength{\emergencystretch}{3em}
\allowdisplaybreaks
\newtheorem*{theorem}{Theorem}
\newtheorem*{thm}{Theorem}
\newtheorem*{lemma}{Lemma}
\newtheorem*{proposition}{Proposition}
\newtheorem*{prop}{Proposition}
\newtheorem*{corollary}{Corollary}
\newtheorem*{cor}{Corollary}
\newtheorem*{claim}{Claim}
\theoremstyle{definition}
\newtheorem*{definition}{Definition}
\newtheorem*{defn}{Definition}
\newtheorem*{remark}{Remark}
\newtheorem*{example}{Example}
\newcommand{\R}{\mathbb R}
\newcommand{\C}{\mathbb C}
\newcommand{\Z}{\mathbb Z}
\newcommand{\Q}{\mathbb Q}
\newcommand{\N}{\mathbb N}
\newcommand{\F}{\mathbb F}
\newcommand{\D}{\mathbb D}
\newcommand{\bB}{\mathbb B}
\newcommand{\bC}{\mathbb C}
\newcommand{\bR}{\mathbb R}
\newcommand{\bZ}{\mathbb Z}
\newcommand{\bN}{\mathbb N}
\newcommand{\bQ}{\mathbb Q}
\newcommand{\bD}{\mathbb D}
\newcommand{\bF}{\mathbb F}
\newcommand{\CP}{\mathbb{CP}}
\newcommand{\RP}{\mathbb{RP}}
\newcommand{\sO}{\mathscr O}
\newcommand{\sI}{\mathscr I}
\newcommand{\sF}{\mathscr F}
\newcommand{\sG}{\mathscr G}
\newcommand{\sH}{\mathscr H}
\newcommand{\sR}{\mathscr R}
\newcommand{\sM}{\mathscr M}
\newcommand{\sV}{\mathscr V}
\newcommand{\sU}{\mathscr U}
\DeclareMathOperator{\Hom}{Hom}
\DeclareMathOperator{\Ext}{Ext}
\DeclareMathOperator{\Tor}{Tor}
\DeclareMathOperator{\Spec}{Spec}
\DeclareMathOperator{\Proj}{Proj}
\DeclareMathOperator{\supp}{supp}
\DeclareMathOperator{\im}{im}
\DeclareMathOperator{\id}{id}
\DeclareMathOperator{\Tr}{Tr}
\DeclareMathOperator{\tr}{tr}
\DeclareMathOperator{\rank}{rank}
\DeclareMathOperator{\ord}{ord}
\DeclareMathOperator{\dist}{dist}
\DeclareMathOperator{\End}{End}
\DeclareMathOperator{\Aut}{Aut}
\DeclareMathOperator{\Gal}{Gal}
\DeclareMathOperator{\vol}{vol}
\DeclareMathOperator{\Cl}{Cl}
\DeclareMathOperator{\coker}{coker}
\DeclareMathOperator{\codim}{codim}
\DeclareMathOperator{\divergence}{div}
\DeclareMathOperator{\Ric}{Ric}
\DeclareMathOperator{\grad}{grad}
\DeclareMathOperator{\Hess}{Hess}
\newcommand{\abs}[1]{\left\lvert#1\right\rvert}
\newcommand{\sabs}[1]{\lvert#1\rvert}
\newcommand{\babs}[1]{\bigl\lvert#1\bigr\rvert}
\newcommand{\Babs}[1]{\Bigl\lvert#1\Bigr\rvert}
\newcommand{\bbabs}[1]{\biggl\lvert#1\biggr\rvert}
\newcommand{\BBabs}[1]{\Biggl\lvert#1\Biggr\rvert}
\newcommand{\norm}[1]{\left\lVert#1\right\rVert}
\newcommand{\snorm}[1]{\lVert#1\rVert}
\newcommand{\bnorm}[1]{\bigl\lVert#1\bigr\rVert}
\newcommand{\Bnorm}[1]{\Bigl\lVert#1\Bigr\rVert}
\newcommand{\bbnorm}[1]{\biggl\lVert#1\biggr\rVert}
\newcommand{\BBnorm}[1]{\Biggl\lVert#1\Biggr\rVert}
\newcommand{\widebar}[1]{\overline{#1}}
\newcommand{\nicefrac}[2]{\frac{#1}{#2}}
\newcommand{\linnprod}[2]{\langle#1,#2\rangle}
\newcommand{\myindex}[1]{#1}
\newcommand{\glsadd}[1]{}
\newcommand{\avoidbreak}{}
\newcommand{\sourceref}[1]{\textnormal{[原文：\nolinkurl{#1}]}}
\newcommand{\thmref}[1]{\sourceref{#1}}
\newcommand{\lemmaref}[1]{\sourceref{#1}}
\newcommand{\propref}[1]{\sourceref{#1}}
\newcommand{\corref}[1]{\sourceref{#1}}
\newcommand{\defnref}[1]{\sourceref{#1}}
\newcommand{\exampleref}[1]{\sourceref{#1}}
\newcommand{\exerciseref}[1]{\sourceref{#1}}
\newcommand{\figureref}[1]{\sourceref{#1}}
\newcommand{\sectionref}[1]{\sourceref{#1}}
\newcommand{\appendixref}[1]{\sourceref{#1}}
\newcommand{\stacksref}[2]{\href{https://stacks.math.columbia.edu/tag/#1}{#2 [#1]}}


\newcommand{\E}{\mathbb E}
\newcommand{\Prob}{\mathbb P}
\newcommand{\Reff}{\mathcal R_{\mathrm{eff}}}
\newcommand{\eps}{\varepsilon}
\DeclareMathOperator{\diam}{diam}
\DeclareMathOperator{\Var}{Var}
\DeclareMathOperator{\Crit}{Crit}
\newcommand{\dd}{\,\mathrm d}

\begin{document}
\section*{072\quad 有限域张量解析秩的精确次可加性}
\subsection*{来源与计数目标}
Shachar Lovett, \emph{The Analytic Rank of Tensors and Its Applications}, Discrete Analysis 2019:7, DOI 10.19086/da.8654。采用期刊所链 arXiv:1806.09179v5（2019-06-03）的 Theorem 1.5、Section 1.2 定义与 Section 2 全部证明，获取2026-09-28。
\par\url{https://arxiv.org/html/1806.09179v5}
\par\url{https://discreteanalysisjournal.com/article/8654-the-analytic-rank-of-tensors-and-its-applications}
\par\url{https://discreteanalysisjournal.com/for-authors}
\par\textbf{本文件只计一个问题（整理者概述）：}在有限域上、$d\geq2$、有限维 $V$ 的范围内证明 $\operatorname{arank}(T+S)\leq\operatorname{arank}(T)+\operatorname{arank}(S)$，不只证明带 $2^d$ 因子的近似版本。
\subsection*{作者命题与人类证明}

\paragraph{Notation from Section 1.2.}
Let $\mathbb F$ be a finite field, and let $\chi:\mathbb F\to\mathbb C$ be a nontrivial additive character. Given a function $F:X\to\mathbb F$, its bias is
\[
\operatorname{bias}(T):=\mathbb E_{x\in X}[\chi(F(x))].
\]
In particular, let $T:V^d\to\mathbb F$ be an order-$d$ tensor. The bias of $T$ is always real and in $(0,1]$. To see that, define $T(\cdot,x^2,\ldots,x^d)$ to be the order-$1$ tensor on $x^1$ given a fixing of $x^2,\ldots,x^d$. Then
\begin{align*}
\operatorname{bias}(T)&=\mathbb E_{x^2,\ldots,x^d\in V}\left[\mathbb E_{x^1\in V}[\chi(T(x^1,\ldots,x^d))]\right]\\
&=\Pr_{x^2,\ldots,x^d\in V}[T(\cdot,x^2,\ldots,x^d)\equiv0]. \tag{1}
\end{align*}
This is since for an order-$1$ tensor (namely, a linear form), its bias is $1$ if it is identically zero, and is $0$ otherwise. The analytic rank of $T$ is defined to be
\[
\operatorname{arank}(T):=-\log_{|\mathbb F|}\operatorname{bias}(T).
\]
\begin{theorem}[1.5]
Let $T,S:V^d\to\mathbb F$ be order-$d$ tensors. Then
\[
\operatorname{arank}(T+S)\leq\operatorname{arank}(T)+\operatorname{arank}(S).
\]
\end{theorem}
\paragraph{2. Proof of Theorem 1.5.}
We prove Theorem 1.5 in this section. It suffices to prove that for any two order-$d$ tensors $T,S:V^d\to\mathbb F$ it holds that
\[
\operatorname{bias}(T+S)\geq\operatorname{bias}(T)\operatorname{bias}(S). \tag{3}
\]
We first introduce some notation. Let $\mathbf x=(x^1,\ldots,x^d),\mathbf y=(y^1,\ldots,y^d)\in V^d$. For $I\subseteq[d]$ define $I^c=[d]\setminus I$ and shorthand $\mathbf x^I:=(x^i:i\in I)$. Define $T_I(\mathbf x,\mathbf y)$ as
\[
T_I(\mathbf x,\mathbf y):=T_I(\mathbf x^I,\mathbf y^{I^c})=T(z^1,\ldots,z^d),\qquad
z^i=\begin{cases}x^i&\text{if }i\in I,\\y^i&\text{if }i\notin I.\end{cases}
\]
That is, $T_I$ is the tensor $T$ evaluated over $\mathbf x^I,\mathbf y^{I^c}$. Observe that $T(\mathbf x+\mathbf y)$ decomposes as the sum
\[
T(\mathbf x+\mathbf y)=\sum_{I\subseteq[d]}T_I(\mathbf x^I,\mathbf y^{I^c}). \tag{4}
\]
Express $\operatorname{bias}(T)\operatorname{bias}(S)$ as
\begin{align*}
\operatorname{bias}(T)\operatorname{bias}(S)
&=\operatorname{bias}(T(\mathbf x)+S(\mathbf y))\\
&=\operatorname{bias}(T(\mathbf x)+S(\mathbf x+\mathbf y)).
\end{align*}
Here, we used the fact that the joint distributions of $(\mathbf x,\mathbf y)$ and $(\mathbf x,\mathbf x+\mathbf y)$ are identical. Next, decompose $S(\mathbf x+\mathbf y)$ using Equation 4 as
\begin{align*}
\operatorname{bias}(T)\operatorname{bias}(S)
&=\operatorname{bias}\left(T(\mathbf x)+\sum_{I\subseteq[d]}S_I(\mathbf x^I,\mathbf y^{I^c})\right)\\
&=\operatorname{bias}\left((T+S)(\mathbf x)+\sum_{I\subsetneq[d]}S_I(\mathbf x^I,\mathbf y^{I^c})\right).
\end{align*}
Fix a choice of $\mathbf y=\mathbf b\in V^d$ so that
\[
\operatorname{bias}(T)\operatorname{bias}(S)\leq
\left|\operatorname{bias}\left((T+S)(\mathbf x)+\sum_{I\subsetneq[d]}S_I(\mathbf x^I,\mathbf b^{I^c})\right)\right|.
\]
Observe that $S_I(\mathbf x^I,\mathbf b^{I^c})$ is an order-$|I|$ tensor of the inputs $\mathbf x^I$. The proof of Theorem 1.5 follows from the following lemma, applied to $R_{[d]}(\mathbf x)=(T+S)(\mathbf x)$ and $R_I(\mathbf x^I)=S_I(\mathbf x^I,\mathbf b^{I^c})$ for $I\subsetneq[d]$.
\begin{lemma}[2.1]
For each $I\subseteq[d]$, let $R_I:V^I\to\mathbb F$ be an order-$|I|$ tensor. Consider the function
\[
R(\mathbf x)=\sum_{I\subseteq[d]}R_I(\mathbf x^I).
\]
Then $|\operatorname{bias}(R)|\leq\operatorname{bias}(R_{[d]})$.
\end{lemma}
In order to prove Lemma 2.1, we first prove the following claim.
\begin{proposition}[Claim 2.2]
Let $W_0,\ldots,W_n:\mathbb F^m\to\mathbb F$ be functions. Consider functions $A,B:\mathbb F^n\times\mathbb F^m\to\mathbb F$ defined as follows:
\[
A(x,y)=\sum_{i=1}^n x_iW_i(y),\qquad B(x,y)=A(x,y)+W_0(y).
\]
Then $|\operatorname{bias}(B)|\leq\operatorname{bias}(A)$.
\end{proposition}
\begin{proof}
We have
\begin{align*}
\operatorname{bias}(B)&=\mathbb E_{x,y}[\chi(B(x,y))]\\
&=\mathbb E_y[1_{W_1(y)=\ldots=W_n(y)=0}\cdot\chi(W_0(y))].
\end{align*}
Hence
\[
|\operatorname{bias}(B)|\leq\mathbb E_y[1_{W_1(y)=\ldots=W_n(y)=0}]=\operatorname{bias}(A).
\]
\end{proof}
\begin{proof}[Proof of Lemma 2.1]
Fix $i\in[d]$ and decompose $R(\mathbf x)$ as
\[
R(\mathbf x)=\sum_{I\subseteq[d],i\in I}R_I(\mathbf x^I)+\sum_{I\subseteq[d],i\notin I}R_I(\mathbf x^I).
\]
Setting $x=x^i$ and $y=\mathbf x^{[d]\setminus\{i\}}$, the first sum has the form $\sum x_iW_i(y)$, and the second sum which does not depend on $x$ is $W_0(y)$. Claim 2.2 then gives that
\[
|\operatorname{bias}(R)|\leq\operatorname{bias}\left(\sum_{I\subseteq[d],i\in I}R_I(\mathbf x^I)\right).
\]
Applying this iteratively for $i=1,\ldots,d$ completes the proof.
\end{proof}

\subsection*{整理说明与先修范围（非作者正文）}
使用范围 $d\geq2$ 在来源栏就地补明，作者定理句保持原样；此范围确保正文中的偏差为正。记号段按原文保留将一般函数 $F$ 的偏差写成 $\operatorname{bias}(T)$ 的字样。只将公式分行和网页数学恢复为独立 TeX。Section 2 的 Claim 2.2 与 Lemma 2.1 是核心支撑，已全录；不收入 Section 3 的其他应用，也不把支撑引理另算题。标准先修为有限域加法特征标正交性、有限平均和多线性展开。难点在对角平移后分离最高阶张量、固定辅助变量并逐坐标消去低阶项，因而保留精确常数1。数学正文据完整HTML转录；PDF截图通道未成功，未标作原PDF下载件。
\subsection*{可修改的简短标注（非作者正文）}
$c$：张量解析秩的加法控制

$a$：加法特征标；偏差与零纤维概率；多线性展开；均匀分布平移；逐坐标条件平均

\texttt{cross\_theory\_status = yes}。把张量求和后的秩不等式转为偏差乘法不等式，以随机平移展开和逐坐标正交性逐次消去低阶项，再取对数返回秩。位置：Section 2 的式(3)--(4)、Claim 2.2 和 Lemma 2.1。

全部标注为非穷尽、可修改初稿。
\subsection*{来源许可与修改记录}
期刊公开作者政策规定刊发文章按 CC BY 许可复用，作者保留版权；本文采用期刊链接的正式刊发版本。arXiv页面自身的存档许可与期刊内容的复用许可分开记录；不将前者误称CC许可。保留作者、论文题名和期刊DOI。
2026-09-28：ChatGPT 整理为独立 TeX，加入编号、来源定位、整理说明与初稿标注；原作者数学论证与整理者文字分开标示。
\end{document}
